\chapter{Systèmes de frontière, rang dimensionnel et fonctionnelle d'action}
\label{chap:pantalla}

En mécanique dimensionnelle, la dimension s'obtient à partir du nombre de canaux
indépendants disponibles à une frontière. Elle n'est pas ajoutée comme étiquette extérieure
à l'espace des états : elle se calcule comme le rang d'un module de prolongements
soumis à ses relations. Un changement de dimension correspond alors à une
variation de rang au cours du transport.

\section{Systèmes de restriction et données de frontière}

Soit
\[
 \cdots\longrightarrow\mathcal S_{n+1}
 \xrightarrow{r_{n+1,n}}\mathcal S_n\longrightarrow\cdots
\]
un système projectif d'états survivants. La structure de frontière au
niveau \(n\) est la paire
\[
 \mathcal F_n=(r_{n+1,n},\beta_n),
\]
où \(\beta_n\) conserve la donnée minimale nécessaire pour décider de la
compatibilité d'un prolongement. Ses coordonnées peuvent inclure le type
arithmétique, l'orientation, le cocycle de retenue et la classe d'incidence.

\begin{lemma}[Compatibilité des restrictions successives]
Supposons
\[
 r_{n+2,n}=r_{n+1,n}\circ r_{n+2,n+1},
 \qquad
 \beta_n\circ r_{n+1,n}
 =b_{n+1,n}\circ\beta_{n+1}.
\]
Alors toute condition de compatibilité exprimée au moyen de \(\beta_{n+1}\)
induit une condition compatible au niveau \(n\).
\end{lemma}
\begin{proof}
Pour \(x\in\mathcal S_{n+2}\), l'égalité des compositions identifie la
restriction directe aux deux restrictions consécutives. La seconde
égalité transporte la donnée de frontière à travers le même diagramme. La
condition au niveau \(n\) est donc l'image de la condition plus
profonde.
\end{proof}

\section{Module de canaux et rang dimensionnel}

Un corps \(\Kfield_n\) étant fixé, soit \(M_n(x)\) le module engendré par les
canaux de prolongement de \(x\), et soit \(\mathcal R_n(x)\) le sous-module des
relations entre eux.

\begin{definition}[Écran dimensionnel]
L'écran associé à \(x\) au niveau \(n\) est
\[
 \Sigma_n(x)=
 (\Kfield_n,M_n(x),\beta_n(x),\mathcal R_n(x)),
\]
et sa dimension effective est
\[
 d_{\Sigma_n}(x)=
 \dim_{\Kfield_n}\frac{M_n(x)}{\mathcal R_n(x)}.
\]
Si \(M_n(x)\cong\Kfield_n^{g_n}\) et si les colonnes d'une matrice
\(R_n(x)\) engendrent \(\mathcal R_n(x)\), alors
\[
 d_{\Sigma_n}(x)=g_n-\rank_{\Kfield_n}R_n(x).
\]
Lorsqu'un supplémentaire est fixé et que \(P_n(x)\) désigne le projecteur dont l'image est
isomorphe au quotient, la même dimension est
\(d_{\Sigma_n}(x)=\rank_{\Kfield_n}P_n(x)\).
\end{definition}

\begin{proposition}[Invariance par équivalence de présentations]
La dimension d'écran est invariante par changements de base inversibles dans
les générateurs et dans les relations.
\end{proposition}
\begin{proof}
Un changement de bases transforme la matrice des relations \(R\) en \(URV\), avec
\(U\) et \(V\) inversibles. Comme \(\rank(URV)=\rank R\), le nombre
\(g-\rank R\) ne change pas. De manière équivalente, deux projecteurs qui présentent le
même quotient ont des images isomorphes et donc le même rang.
\end{proof}

Pour un transport \(T_\gamma:\mathcal S_n\to\mathcal S_m\), nous définissons
\[
 \Delta d(\gamma;x)=
 d_{\Sigma_m}(T_\gamma x)-d_{\Sigma_n}(x).
\]
Une transition dimensionnelle transporte simultanément l'état et son module
de canaux ; la dynamique à rang fixe est le cas \(\Delta d=0\).

\section{Écrans arithmétiques et réticulaires hérités de HMT}

Les rangs qui interviennent plus loin ne sont pas des étiquettes ajoutées aux
équations. Ils proviennent de constructions linéaires distinctes et conservent leurs
applications de sélection. Le code ternaire de Golay poinçonné
\[
 G_{11}=[11,6,5]_3\subset\mathbb F_3^{11}
\]
est parfait de rayon deux \cite{MacWilliamsSloane1977}, car les boules
centrées sur ses \(3^6\) mots sont
disjointes et
\[
 V_3(11,2)=1+2\binom{11}{1}+2^2\binom{11}{2}
 =1+22+220=243,
\qquad
 3^6\cdot243=3^{11}.
\]
Ici, \(243\) est le cardinal de l'espace des syndromes. Il ne s'identifie pas par
simple égalité numérique à la couche de \(243\) cellules de la complétion
cubique d'APP\@.

La seconde chaîne part du module de permutation
\(V_{12}=\mathbb R^{12}\) de \(A_5\) sur les douze sommets de l'icosaèdre.
Pour fixer la jauge, nous faisons agir \(A_5\) sur
\(\{0,1,2,3,4\}\), prenons
\(H=\langle(0\ 1\ 2\ 3\ 4)\rangle\cong C_5\) et identifions la position
\(j\) à la classe à gauche \(r_jH\). En notation des images, les douze
représentants sont
\[
\begin{array}{c|c@{\qquad}c|c@{\qquad}c|c}
j&r_j&j&r_j&j&r_j\\ \hline
1&(3,2,4,1,0)&5&(3,2,0,4,1)&9&(3,1,2,4,0)\\
2&(3,1,0,2,4)&6&(4,0,1,2,3)&10&(2,1,4,3,0)\\
3&(0,1,3,4,2)&7&(4,3,2,1,0)&11&(3,4,1,2,0)\\
4&(1,4,2,3,0)&8&(0,2,3,1,4)&12&(4,2,1,3,0).
\end{array}
\]
Nous notons \(5A\) la classe de conjugaison contenant le générateur
\(g_5=(0\ 1\ 2\ 3\ 4)\) de \(H\), et \(5B\) la classe contenant
\(g_5^2\). Dans l'ordre des classes \(1A,2A,3A,5A,5B\), le caractère de permutation est
\(\chi_{12}=(12,0,0,2,2)\). Pour
\(\phi=(1+\sqrt5)/2\) et \(\phi'=(1-\sqrt5)/2\), nous choisissons le caractère
irréductible
\[
 \chi_3=(3,-1,0,\phi,\phi').
\]
Le produit scalaire des caractères donne la décomposition exacte
\[
 V_{12}\cong V_1\oplus V_3\oplus V_{3'}\oplus V_5.
\]

Le projecteur
\[
 P_{11}=I_{12}-\frac1{12}J_{12}
\]
élimine exclusivement le vecteur uniforme et a pour image
\[
 V_{11}=\mathbf1^\perp\cong V_3\oplus V_{3'}\oplus V_5.
\]
Pour l'état dodécaphasique
\[
 K=(234,543,140,729,659,824,621,58,914,794,146,601),
 \qquad \sum_jK_j=6263,
\]
soit \(\widetilde K=P_{11}K\). La correspondance position--classe précédente
définit les matrices de permutation \(\rho(g)\), et le projecteur central
\[
 P_3=\frac3{60}\sum_{g\in A_5}\chi_3(g^{-1})\rho(g)
\]
satisfait \(P_3^2=P_3=P_3^{\mathsf T}\) et a pour rang trois. L'évaluation
exacte dans \(\mathbb Q(\sqrt5)\) produit
\[
 u_K=P_3\widetilde K,
 \qquad
 \|u_K\|^2=\frac{6638585+2275584\sqrt5}{20}>0.
\]
Par conséquent, la droite \(\ell_K=\mathbb Ru_K\) est bien définie et
\[
 V_{10}:=(V_3\cap\ell_K^\perp)\oplus V_{3'}\oplus V_5,
 \qquad
 V_{11}=\ell_K\oplus V_{10},
 \qquad \dim V_{10}=10.
\]
La chaîne \(12\to11\to10\) signifie, respectivement, éliminer la
composante uniforme et la droite polarisée par \(K\). L'ensemble fini des
soixante éléments se divise en douze classes ; son énumération
reconstruit exactement le projecteur et la norme.

La réduction \(26\to24\) appartient à une autre catégorie. Soit \(U\) le plan
hyperbolique entier de matrice de Gram
\(\left(\begin{smallmatrix}0&1\\1&0\end{smallmatrix}\right)\), y sea
\[
 L_{26}=\Lambda_{24}\oplus U.
\]
Le facteur hyperbolique étant fixé, il s'agit de la présentation standard du
réseau pair unimodulaire lorentzien \(II_{25,1}\)
\cite{ConwaySloane1999}.
Pour un vecteur isotrope primitif \(e_+\in U\),
\[
 e_+^\perp=\Lambda_{24}\oplus\mathbb Ze_+,
\qquad
 e_+^\perp/\mathbb Ze_+\cong\Lambda_{24}.
\]
La diminution du rang de deux provient de l'imposition de l'orthogonalité, puis du passage au quotient
par la droite nulle. Ce n'est pas la réduction \(11\to10\) répétée avec d'autres
nombres.

\section{Involution entre donnée résiduelle et quotient euclidien}

La décomposition \(n=9Q+r\) suggère deux coordonnées de nature différente :
le représentant résiduel \(r\) et le quotient \(Q\). Pour formuler leur
échange dans un domaine stable, nous considérons
\[
 \mathcal D=\mathbb Z^2\times\mathbb R_{>0},
\qquad
 E_R(m,w)=\frac{m^2}{R^2}+w^2R^2,
\]
et nous définissons
\[
 \mathsf T(m,w,R)=(w,m,R^{-1}).
\]
\begin{proposition}[Dualité résiduelle--euclidienne]
\(\mathsf T\) est une involution et
\[
 E_R(m,w)=E_{R^{-1}}(w,m).
\]
Sur \(\ell^2(\mathbb Z^2)\), l'unitaire
\(U_T|m,w\rangle=|w,m\rangle\) satisfait
\[
 U_TH(R)U_T^{-1}=H(R^{-1}),
\qquad
 H(R)|m,w\rangle=E_R(m,w)|m,w\rangle.
\]
\end{proposition}
\begin{proof}
La seconde application échange de nouveau les coordonnées et inverse deux
fois \(R\). L'identité d'énergie s'obtient par substitution directe ; l'
identité opératorielle se vérifie sur la base orthonormée.
\end{proof}

La proposition est une dualité réticulaire exacte. La lecture physique en termes d'
impulsion, d'enroulement et de rayon de compactification requiert une application supplémentaire
vers l'espace des états d'une théorie des cordes ; elle ne se déduit pas de l'
involution algébrique \cite{Giveon1994}.

\section{Interface typée avec la réduction undécadimensionnelle}

L'écran \(V_{11}=\ell_K\oplus V_{10}\) et l'involution
\((m,w,R)\mapsto(w,m,R^{-1})\) fournissent les deux données algébriques qui
motivent la comparaison avec le contexte de la théorie M : une réduction marquée
de onze à dix coordonnées et une équivalence de rayons réciproques. Pour que
cette comparaison soit une réalisation physique, il faut spécifier un morphisme
\[
 \mathcal R_M:
 (\mathcal X_{\rm TPK},V_{11},\ell_K,\mathsf T)
 \longrightarrow
 (\mathcal H_M,\mathcal F_{10},R_{\rm phys},\alpha',
 \mathcal Q_{\rm susy}),
\]
où le codomaine comprend un espace des états, des champs
décadimensionnels, un rayon physique, une tension de corde et des charges supersymétriques
\cite{Witten1995}. La décomposition algébrique fixe le domaine de
\(\mathcal R_M\) ; elle n'introduit pas à elle seule des oscillateurs, des branes,
la supersymétrie ni les équations de supergravité.

Cette séparation conserve la profondeur de la correspondance sans altérer
les types. La réduction \(11\to10\), l'involution des rayons et le quotient
\(26\to24\) sont des théorèmes des objets définis. L'affirmation qu'une
théorie dynamique précise en est l'image exige d'exhiber \(\mathcal R_M\) et de prouver
qu'il entrelace les actions, les restrictions et les observables.

\section{État dimensionnel et transport simultané}

L'état sur lequel agit la mécanique dimensionnelle s'écrit
\[
 \mathfrak x=(x,\Sigma_n(x),\beta_n(x),[\gamma]),
\]
où \(x\in\mathcal X_{\rm TPK}\), \(\Sigma_n(x)\) est son écran,
\(\beta_n(x)\) est la donnée de frontière et \([\gamma]\) est la classe de trajectoire
compatible avec cette donnée. Un transport n'agit pas seulement sur \(x\) : il se compose de
\[
 (T_\gamma,F_\gamma):
 (x,\Sigma_n,\beta_n)\longrightarrow
 (T_\gamma x,\Sigma_m,\beta_m),
\]
où \(F_\gamma:M_n(x)\to M_m(T_\gamma x)\) envoie les relations sur des relations
et satisfait le diagramme de compatibilité avec \(\beta_n,\beta_m\). Le saut
dimensionnel est le changement de rang du morphisme induit
\[
 \overline F_\gamma:
 M_n(x)/\mathcal R_n(x)\longrightarrow
 M_m(T_\gamma x)/\mathcal R_m(T_\gamma x).
\]
Cette formulation évite d'attribuer une dimension au point isolé : la dimension
appartient au système de canaux qui accompagne l'état et son transport.

\section{Descente vers une dynamique sur des états ponctuels}

Soit \(q_{\rm pt}:\mathcal X_{\rm TPK}\to Z\) une projection qui élimine la
généalogie de l'état, et définissons
\[
 Y_{\rm pt}:=\operatorname{im}(q_{\rm pt}).
\]
Soit, en outre, \(U:\mathcal X_{\rm TPK}\to\mathcal X_{\rm TPK}\) une
évolution complète.

\begin{theorem}[Critère de descente par fibres]
Il existe une unique application
\(\overline U:Y_{\rm pt}\to Y_{\rm pt}\) telle que
\(q_{\rm pt}U=\overline Uq_{\rm pt}\) si et seulement si
\[
 q_{\rm pt}(x)=q_{\rm pt}(y)
 \Longrightarrow q_{\rm pt}(Ux)=q_{\rm pt}(Uy).
\]
\end{theorem}
\begin{proof}
La nécessité résulte de l'application de \(\overline U\) à deux représentants d'une
même fibre. Pour la suffisance, on définit
\[
 \overline U(q_{\rm pt}(x))=q_{\rm pt}(Ux).
\]
La condition garantit que la définition ne dépend pas du représentant. Comme
le codomaine a été fixé égal à l'image de \(q_{\rm pt}\), la
surjectivité requise pour l'unicité est automatique.
\end{proof}

Lorsque la condition échoue, une même classe ponctuelle possède des images distinctes ;
l'évolution descendue est alors multivoque. Conserver la donnée de
mémoire, restreindre le domaine ou travailler avec une correspondance sont trois
modifications mathématiques différentes.

\section{Fonctionnelle additive sur les trajectoires orientées}

Pour éviter un terme caché de composition des trajectoires, l'état d'une arête est
étendu par la direction et le type arithmétique de l'arête précédente. Si
\(\widetilde e_j=(e_j,d_{j-1},t_{j-1})\), nous définissons les poids locaux
\[
 w_{\rm giro}(\widetilde e_j)=\mathbf1_{d_j\ne d_{j-1}},
 \qquad
 w_{\rm tipo}(\widetilde e_j)=\mathbf1_{t_j\ne t_{j-1}}.
\]
Les données \((d_0,t_0)\) font partie de la condition initiale de frontière.
Pour une trajectoire augmentée
\(\widetilde\gamma=(\widetilde e_1,\ldots,\widetilde e_R)\), soient
\[
 N_{\rm giro}(\widetilde\gamma)=\sum_{j=1}^R
 w_{\rm giro}(\widetilde e_j),
 \qquad
 N_{\rm tipo}(\widetilde\gamma)=\sum_{j=1}^R
 w_{\rm tipo}(\widetilde e_j).
\]
On définit
\[
 \mathbf I(\widetilde\gamma)=
 (N_{\rm giro}(\widetilde\gamma),N_{\rm tipo}(\widetilde\gamma)).
\]
La composition n'est permise que lorsque la donnée précédente du second chemin
coïncide avec la donnée terminale du premier. Avec cette convention, le partage
de la somme entre les deux intervalles d'arêtes démontre
\[
 \mathbf I(\widetilde\gamma\widetilde\eta)
 =\mathbf I(\widetilde\gamma)+\mathbf I(\widetilde\eta),
\]
de sorte que \(\mathbf I\) est un cocycle à valeurs dans \(\NN^2\). Pour
\(\lambda>0\), la combinaison
\[
 I_\lambda(\widetilde\gamma)
 =N_{\rm giro}(\widetilde\gamma)
 +\lambda N_{\rm tipo}(\widetilde\gamma)
\]
est une fonctionnelle scalaire sur l'espace des trajectoires.

Dans la formulation lagrangienne ordinaire, la fonctionnelle d'action est définie
sur des courbes d'un espace de configuration choisi au préalable. Ici, l'
ordre est antérieur : \(I_\lambda\) quantifie les changements de direction et de loi
arithmétique dans la catégorie même des trajectoires augmentées. Une application de
réalisation peut ensuite transporter ce cocycle vers une action
continue. L'additivité démontrée ci-dessus est la propriété qui permet cette
comparaison sans identifier d'avance les deux fonctionnelles.

\begin{principle}[Interaction minimale]
Parmi les trajectoires admissibles dont les extrémités et les données de frontière sont fixées, l'
évolution sélectionnée minimise \(I_\lambda\).
\end{principle}

\begin{proposition}[Existence finie de minimiseurs]
Si \(\Gamma(a,b)\) est fini et non vide, \(I_\lambda\) atteint un minimum sur
\(\Gamma(a,b)\).
\end{proposition}
\begin{proof}
L'image d'un ensemble fini non vide par une fonction réelle possède un
élément minimal.
\end{proof}

Le caractère de masse du \cref{chap:masas} constitue une autre fonctionnelle sur les
trajectoires : il transforme un cocycle entier en un rapport positif. La
fonctionnelle et le mécanisme qui sélectionne la trajectoire sont des objets distincts ;
leur composition sera spécifiée avant l'introduction d'une unité de
masse.
